\section*{2026 年 3 月 2 日作业}
\phantomsection
\addcontentsline{toc}{section}{2026 年 3 月 2 日作业}
\noindent\textbf{截止时间：2026 年 3 月 9 日 13:30}

\begin{problem}
设 $m$ 和 $n$ 为给定的正整数，$\omega = \exp(2\pi \ii / m)$。证明：对任何 $A, B \in \C^{n\times n}$ 都有
\[
A^m + B^m = \frac{1}{m}\sum_{k=0}^{m-1} \left(A + \omega^k B\right)^m.
\]
注：$\exp(\ii\theta) = \ee^{\ii\theta} = \cos\theta + \ii\sin\theta$。

\begin{solution}
% TODO: 在此填写解答。
\end{solution}
\end{problem}

\begin{problem}
设 $n$ 为给定的正整数，$\omega = \exp(2\pi \ii / n)$，$A \in \C^{n\times n}$。若 $A$ 的 $(i,j)$ 元素为
\[
A_{ij} = \omega^{(i-1)(j-1)},
\]
求 $A^{-1}$。

\begin{solution}
% TODO: 在此填写解答。
\end{solution}
\end{problem}

\begin{problem}
证明：复数 $z_1, z_2, z_3$ 满足
\[
|z_1-z_2| = |z_2-z_3| = |z_3-z_1|
\]
的充要条件是
\[
z_1^2 + z_2^2 + z_3^2 - z_1z_2 - z_2z_3 - z_1z_3 = 0.
\]
再利用此工具证明 Napoleon 定理：在欧氏平面内，由三角形 $ABC$ 的各边分别向外作正三角形，则这三个正三角形的中心两两相连可构成正三角形。

\begin{solution}
% TODO: 在此填写解答。
\end{solution}
\end{problem}

\begin{problem}
已知
\[
A=
\begin{bmatrix}
-1 & 4 & -3 & 1 \\
-8 & 23 & -21 & 7 \\
-26 & 74 & -70 & 23 \\
-61 & 174 & -166 & 54
\end{bmatrix}.
\]
求一个非奇异矩阵 $X$ 和一个上三角矩阵 $T$，使得
\[
A = XTX^{-1}.
\]

\begin{solution}
% TODO: 在此填写解答。
\end{solution}
\end{problem}

\begin{problem}
计算
\[
\oint_{|z|=1} \overline{z}\,\dd z.
\]
此处的积分路径沿单位圆逆时针方向。

\begin{solution}
% TODO: 在此填写解答。
\end{solution}
\end{problem}

\begin{problem}[选做]
解方程：
\[
8x^3 - 6x = 1.
\]

\begin{solution}
% TODO: 在此填写解答。
\end{solution}
\end{problem}

\begin{problem}[选做]
定义
\[
\mathbb{H} = \left\{
\begin{bmatrix}
\alpha & \beta \\
-\overline{\beta} & \overline{\alpha}
\end{bmatrix}
: \alpha,\beta\in\C
\right\}.
\]
证明 $\mathbb{H}$ 对加法和乘法封闭，并建立这两种运算所满足的运算律。

\begin{solution}
% TODO: 在此填写解答。
\end{solution}
\end{problem}

\begin{problem}[选做]
给定实数 $\alpha,\beta$。证明：
\[
\left\{\alpha\cos\theta + \beta\sin\theta\cdot \ee^{\ii\phi} : \theta,\phi\in[0,2\pi]\right\}
=
\left\{\xi\in\C : \frac{(\operatorname{Re}\xi)^2}{\alpha^2+\beta^2} + \frac{(\operatorname{Im}\xi)^2}{\beta^2} \le 1\right\}.
\]

\begin{solution}
% TODO: 在此填写解答。
\end{solution}
\end{problem}
